%% uniform_rank_bounds_2026-08-05.tex
%%
%% Machine-checked canonical heights and rank lower bounds, uniformly in
%% the curve.
%%
%% Sequel to Papers/formal_mordell_weil_rank_2026-07-30.tex, cited here
%% as [MW].
%%
%% Every mathematical claim below is either a kernel-verified Lean 4
%% theorem at HEAD b99f2ab6 (2026-08-05), axioms a subset of
%% [propext, Classical.choice, Quot.sound], zero project axioms, zero
%% sorries, no native_decide, or standard cited literature.  Verified
%% build data: lake build PF = 4711 jobs, mathlib pin eed770a4,
%% toolchain leanprover/lean4:v4.24.0-rc1.

\documentclass[11pt,reqno]{amsart}

\usepackage[T1]{fontenc}
\usepackage{lmodern}
\usepackage{amsmath,amssymb,amsthm}
\usepackage{mathtools}
\usepackage{microtype}
\usepackage{geometry}
\geometry{margin=1in}
\usepackage[hidelinks]{hyperref}
\usepackage{enumitem}
\usepackage{booktabs}
\usepackage{longtable}
\usepackage{seqsplit}

\emergencystretch=3em
\hbadness=6000
\sloppy

\theoremstyle{plain}
\newtheorem{theorem}{Theorem}[section]
\newtheorem{lemma}[theorem]{Lemma}
\newtheorem{proposition}[theorem]{Proposition}
\newtheorem{corollary}[theorem]{Corollary}

\theoremstyle{definition}
\newtheorem{definition}[theorem]{Definition}

\theoremstyle{remark}
\newtheorem{remark}[theorem]{Remark}

% \lt: Lean identifiers and shell commands.  Plain \texttt, NOT \seqsplit --
% seqsplit deletes spaces (it would print "lake build PF" as "lakebuildPF")
% and permits a line break after any single character.
\DeclareRobustCommand{\lt}[1]{\texttt{#1}}
% \ltw: long unbroken digit strings, where breaking anywhere is what we want.
\DeclareRobustCommand{\ltw}[1]{\texttt{\seqsplit{#1}}}
\newcommand{\hh}{\widehat{h}}
\newcommand{\nh}{H}
\DeclareMathOperator{\rank}{rank}
\DeclareMathOperator{\reg}{Reg}
\DeclareMathOperator{\Res}{Res}
\newcommand{\DDz}{\mathcal{D}}
\newcommand{\Sfz}{\mathcal{S}}
\newcommand{\Prz}{\mathcal{P}}

\title[Machine-checked canonical heights and rank lower bounds,
uniformly in the curve]{Machine-checked canonical heights and rank lower
bounds, uniformly in the curve}

\author{Pablo Cohen}
\address{Mesa, Arizona, USA}
\email{psolorzano@alumni.berklee.edu}

\date{5 August 2026}

\subjclass[2020]{Primary 11G05, 11Y16; Secondary 68V20, 68V15}

\keywords{elliptic curve, canonical height, Mordell--Weil rank,
independence of points, resultant, Lean 4, interactive theorem proving}

\begin{document}

\begin{abstract}
We report a machine-checked construction of the canonical
(N\'eron--Tate) height, and of rank lower bounds, for elliptic curves
over $\mathbb{Q}$, which is \emph{uniform in the curve}. Working in
Lean~4 over \texttt{mathlib}: for every Weierstrass curve $W/\mathbb{Q}$
whose $b$-invariants are integers, whose discriminant $\Delta$ is
nonzero, and which has no rational $2$-torsion, we construct
$\hh \colon W(\mathbb{Q}) \to \mathbb{R}$, prove the parallelogram law
for it \emph{exactly}, deduce $\hh(kR) = k^2\hh(R)$ and bi-additivity of
the associated pairing, and conclude: if $n$ rational points have
nonzero Gram determinant, then $n \le \rank_{\mathbb{Z}}
W(\mathbb{Q})$. Every constant is an explicit polynomial in the
coefficients of the curve. The predecessor of this paper proved
$\rank \ge 2$ for \texttt{389a1} and $\rank \ge 3$ for \texttt{5077a1}
through certificate suites built by hand for each curve; here both are
re-derived by instantiating the uniform machine, with the earlier
interval arithmetic as the only per-curve numerical input.

The mathematical content is a reduction. The secant forms that control
$x(P{+}Q)$ and $x(P{-}Q)$ restrict on the diagonal to the duplication
pair $(\Phi,\Psi)$, so the secant content problem reduces to the
resultant identity $\Res(\Phi,\Psi) = \Delta^2$ rather than needing an
elimination theory of its own; twin B\'ezout certificates in the two
monomial weights, together with a cubing of the duplication identity,
give a divisor bound $\Delta^6$ for the content of the secant triple, on
every curve. Exactly one hypothesis survives universalization --- no
rational $2$-torsion --- and it is necessary, because $\Psi$ vanishes
precisely at rational $2$-torsion. The price of uniformity is paid in
the size of the constants, which we tabulate. Kernel axioms are
$[\mathtt{propext}, \mathtt{Classical.choice}, \mathtt{Quot.sound}]$
throughout, with no \lt{sorry}, no \lt{native\_decide} and no
project-level axiom. These are rank \emph{lower} bounds; nothing here
bears on descent, on $L$-functions, or on the Birch--Swinnerton-Dyer
conjecture.
\end{abstract}

\maketitle

\section{Introduction}

\subsection{What is new}

In \cite{MW} the canonical height of N\'eron and Tate was constructed
inside a proof assistant for the first time, and used to prove
$2 \le \rank E(\mathbb{Q})$ for $E = \mathtt{389a1}$ and
$3 \le \rank E'(\mathbb{Q})$ for $E' = \mathtt{5077a1}$. Two components
of that work were already curve-independent: the passage from a doubling
window to a canonical height, and the passage from a nonvanishing Gram
determinant to a rank bound. Everything between them --- the duplication
estimates, the secant estimates, the quasi-parallelogram law and hence
bi-additivity --- was written out by hand, once per curve, in some forty
files.

This paper reports that the whole of the middle is now
curve-independent. The main theorem is the following.

\begin{theorem}[\lt{rank\_ge\_universal}]\label{thm:main}
Let $W$ be a Weierstrass curve over $\mathbb{Q}$ (\texttt{mathlib}'s
\lt{WeierstrassCurve.Affine} $\mathbb{Q}$), and suppose
\begin{enumerate}[label=\textup{(\roman*)}]
  \item the $b$-invariants of $W$ are integers,
        $b_i = B_i$ with $B_2,B_4,B_6,B_8 \in \mathbb{Z}$;
  \item $\Delta := -B_2^2B_8 + 9B_2B_4B_6 - 8B_4^3 - 27B_6^2 \ne 0$;
  \item $W(\mathbb{Q})$ has no rational $2$-torsion: $2R = O$ implies
        $R = O$.
\end{enumerate}
Let $\hh$ be the canonical height of \S\ref{sec:window} and
$\langle\cdot,\cdot\rangle$ the pairing of \S\ref{sec:par}. If
$P \colon \{1,\dots,n\} \to W(\mathbb{Q})$ satisfies
\[
  \det\bigl(\langle P_i, P_j\rangle\bigr)_{i,j} \;\ne\; 0,
\]
then $n \le \rank_{\mathbb{Z}} W(\mathbb{Q})$.
\end{theorem}

We call a curve satisfying (i)--(iii) \emph{admissible}. Nothing else is
assumed: no reduced model, no minimality, no bound on the conductor, no
knowledge of the torsion subgroup beyond (iii). All of the constants
that appear in the construction are explicit polynomials in
$\lvert B_2\rvert, \lvert B_4\rvert, \lvert B_6\rvert, \lvert
B_8\rvert$.

Two specializations are proved in the same file, because they are the
shapes one actually checks. At $n = 2$ the Gram determinant is the
classical regulator expression,
\begin{equation}\label{eq:reg2}
  \det\bigl(\langle P_i,P_j\rangle\bigr)
   = \hh(P)\,\hh(Q) - \langle P,Q\rangle^2
   \qquad (\lt{gram\_det\_fin\_two}),
\end{equation}
so a nonzero regulator forces $\rank \ge 2$
(\lt{rank\_ge\_two\_universal}); at $n = 3$ the corresponding expansion
is \lt{gram\_det\_fin\_three}.

Theorem~\ref{thm:main} rests on the following, which is where the
mathematical work is.

\begin{theorem}[\lt{canheight\_parallelogram\_universal}, \lt{par},
\lt{mul}, \lt{biAdditive\_pairingW}]\label{thm:height}
Let $W$ be admissible. There is a function
$\hh \colon W(\mathbb{Q}) \to \mathbb{R}$, defined as
\[
  \hh(R) \;=\; \lim_{n\to\infty}
     \frac{\log \nh\bigl(x(2^nR)\bigr)}{4^n},
\]
such that, formally:
\begin{enumerate}[label=\textup{(\arabic*)}]
  \item $\hh \ge 0$, $\hh(O) = 0$ and $\hh(-R) = \hh(R)$;
  \item $\bigl\lvert \hh(R) - \log \nh(x(R))\bigr\rvert \le
        \tfrac13 \log \kappa$, where $\kappa$ is the explicit constant
        of \S\ref{sec:dup};
  \item $\hh(P+Q) + \hh(P-Q) = 2\hh(P) + 2\hh(Q)$ for \emph{all}
        $P, Q \in W(\mathbb{Q})$, with no side conditions;
  \item $\hh(kR) = k^2\hh(R)$ for all $k \in \mathbb{Z}$ and all $R$;
  \item the pairing $\langle P,Q\rangle :=
        \tfrac12\bigl(\hh(P{+}Q) - \hh(P) - \hh(Q)\bigr)$ is symmetric
        and additive in each variable, and
        $\langle P,P\rangle = \hh(P)$.
\end{enumerate}
\end{theorem}

Here $\nh(a/b) = \max(\lvert a\rvert, b)$ for a reduced fraction, and
$x(R)$ is the $x$-coordinate of $R$, with $x(O) := 0$ (harmless, since
$\nh(0) = 1$). Every statement in Theorems~\ref{thm:main}
and~\ref{thm:height} is a Lean~4 theorem; \lt{\#print axioms} on each
reports a subset of $[\mathtt{propext}, \mathtt{Classical.choice},
\mathtt{Quot.sound}]$ and nothing else. \S\ref{sec:verif} gives the
theorem-by-theorem map to the sources and the verification data;
\S\ref{sec:notproved} states, specifically, what is \emph{not} proved
--- these are rank lower bounds, and nothing here concerns descent or
$L$-functions.

Note the strengthening in Theorem~\ref{thm:height}(3) and~(4). In
\cite{MW} the parallelogram law carried four non-torsion side
conditions, discharged for the two worked curves by proving their
torsion subgroups trivial. Here it is unconditional on any admissible
curve: the degenerate configurations are closed separately, and the one
place where genuine torsion could intrude --- $2$-torsion --- is
excluded by hypothesis (iii). That is what makes the multiple law and
bi-additivity come out without bookkeeping.

\subsection{Relation to the predecessor, and what the two curves are for}

The results of \cite{MW} are re-derived here, through the uniform
machine:
\begin{align}
  2 &\le \rank_{\mathbb{Z}} E(\mathbb{Q}), \quad
    E = \mathtt{389a1} \colon y^2 + y = x^3+x^2-2x
    &&(\lt{E389a1\_rank\_ge\_two\_universal}), \label{eq:389}\\
  3 &\le \rank_{\mathbb{Z}} E'(\mathbb{Q}), \quad
    E' = \mathtt{5077a1} \colon y^2+y = x^3-7x+6
    &&(\lt{E5077a1\_rank\_ge\_three\_universal}). \label{eq:5077}
\end{align}
These are not new arithmetic facts, and they were not new in \cite{MW}
either: the ranks of these two curves are classical. Their role here is
different. They are the evidence that the uniform machine is not
vacuous and that it reproduces, rather than merely parallels, the
hand-built results: \eqref{eq:389} and \eqref{eq:5077} are obtained by
instantiating Theorem~\ref{thm:main} at the two hypothesis packages and
feeding it the certified interval arithmetic of \cite{MW} --- the
regulator positivity computations --- as the sole numerical input. No
height estimate specific to either curve appears anywhere in the chain.
\S\ref{sec:validation} describes the transfer and records a
verification-methodology point that we think is of independent
interest: the uniform canonical height and the per-curve one turn out to
be equal \emph{by \lt{rfl}}.

\subsection{The obstruction, and how it is avoided}\label{sec:ideas}

The reason the middle of \cite{MW} was per-curve is easy to state. Both
halves of the construction --- doubling and addition --- need two
things about an explicit pair of integer forms: a bound on their
\emph{content} (the common factor that cancels when the fraction is
reduced), and two-sided bounds on their \emph{size}. Per curve, the
content is computed by an extended Euclidean algorithm on the actual
integer polynomials, and the size by summing actual coefficients. Both
outputs are numbers, not formulas, and both computations must be redone
from scratch for each new curve.

The duplication half was made uniform by a resultant identity: writing
$x = X/Y$ and homogenizing,
\begin{equation}\label{eq:PhiPsi}
  \Phi(X,Y) = X^4 - b_4X^2Y^2 - 2b_6XY^3 - b_8Y^4, \qquad
  \Psi(X,Y) = 4X^3Y + b_2X^2Y^2 + 2b_4XY^3 + b_6Y^4,
\end{equation}
one has $\Res(\Phi,\Psi) = \Delta^2$ as an identity in
$\mathbb{Z}[b_2,b_4,b_6,b_8]$ modulo the relation
$4b_8 = b_2b_6 - b_4^2$, and the Sylvester cofactors give two explicit
B\'ezout witnesses. That was already in place, and is summarized in
\S\ref{sec:dup}. The secant half looked as though it would need its own
elimination theory --- a resultant in two sets of variables, on forms
twice as large. It does not, and the reason is the first of four ideas
this paper contributes.

\emph{First: diagonal degeneration.} Write $x_i = u_i/v_i$ in lowest
terms. The three integer forms controlling addition are
$\DDz = (u_1v_2 - u_2v_1)^2$ and the bidegree-$(2,2)$ forms $\Sfz$,
$\Prz$ of \eqref{eq:secforms} below, which satisfy
$\DDz\cdot\bigl(x(P{+}Q) + x(P{-}Q)\bigr) = \Sfz$ and
$\DDz\cdot x(P{+}Q)\,x(P{-}Q) = \Prz$. Restricted to the diagonal
$(u_2,v_2) = (u_1,v_1)$ they \emph{are} the duplication pair:
\begin{equation}\label{eq:diag}
  \Sfz(u,v,u,v) = \Psi(u,v), \qquad
  \Prz(u,v,u,v) = \Phi(u,v), \qquad
  \DDz(u,v,u,v) = 0 .
\end{equation}
The secant construction degenerates, honestly, onto the duplication
construction. Consequently the secant content problem does not need a
new keystone: eliminating the second point against $\DDz$ lands on
$\Phi$ and $\Psi$ at the first point, and $\Res(\Phi,\Psi) = \Delta^2$
finishes the job.

\emph{Second: twin certificates.} Elimination produces the duplication
forms \emph{squared}, weighted by a monomial in the second point's
coordinates. Sylvester's matrix supplies the resultant in both monomial
weights, and both are needed: one gets
\[
  \sigma_1\DDz + \sigma_2\Sfz = v_2^3\,\Psi_1^2
  \qquad\text{and}\qquad
  \sigma_1'\DDz + \sigma_2'\Sfz = u_2^3\,\Psi_1^2,
\]
with explicit integer cofactors, and similarly for $\Prz$ and $\Phi_1$.
Neither weight alone is usable --- $u_2$ and $v_2$ can each be
arbitrarily divisible --- but since $\gcd(u_2,v_2) = 1$, so are
$u_2^3$ and $v_2^3$, and the two weights cancel each other. A common
divisor of the secant triple therefore divides $\Psi_1^2$ and
$\Phi_1^2$ outright.

\emph{Third: the squaring trick.} What is now available are the
\emph{squares} $\Phi_1^2, \Psi_1^2$, while the duplication B\'ezout
identity is linear, $F_1\Phi + G_1\Psi = \Delta^2X^7$. Cubing it puts
every term into the ideal generated by $\Phi^2$ and $\Psi^2$:
\begin{equation}\label{eq:cube}
  (F\Phi + G\Psi)^3
   = \bigl(F^3\Phi + 3F^2G\Psi\bigr)\Phi^2
     + \bigl(3FG^2\Phi + G^3\Psi\bigr)\Psi^2 .
\end{equation}
Hence a common divisor of the secant triple divides
$(\Delta^2X^7)^3 = \Delta^6X^{21}$ and likewise $\Delta^6Y^{21}$, and
coprimality of $X,Y$ leaves $\Delta^6$. That is exactly where the
exponent~$6$ comes from; it is the price of the cube, and we do not
claim it is optimal.

\emph{Fourth: the honest residue.} Per curve, \cite{MW} discharged all
non-torsion side conditions from the triviality of the torsion
subgroup, proved by enumeration. Uniformly, exactly one hypothesis
survives, and it is hypothesis (iii): no rational $2$-torsion. It is not
an artifact. On the curve one has
$\Psi(x) = (2y + a_1x + a_3)^2$ (\lt{psiX\_eq\_sq}), so
$\Psi$ vanishes at $x(R)$ precisely when $2R = O$; at such a point the
duplication map has no denominator, the height window's construction has
nothing to say, and the whole chain stops. Curves with rational
$2$-torsion are genuinely outside the scope of Theorem~\ref{thm:main};
they would need a per-step variant, which we have not carried out.

\subsection{Related formalization}\label{sec:related}

\texttt{mathlib} \cite{mathlib} provides the Weierstrass curve type, the
group law on the affine, projective and Jacobian models with a full
associativity proof, variable change, normal forms, division
polynomials, and reduction over discrete valuation rings. At the commit
pinned for this work (\S\ref{sec:verif}) it contains no arithmetic
height theory, no Mordell--Weil theorem, and no descent; we checked this
directly against the pinned copy.

On \texttt{mathlib} \texttt{master} at the time of writing there is a
\lt{Mathlib/NumberTheory/Height/} directory --- \lt{Basic},
\lt{Northcott}, \lt{MvPolynomial}, \lt{NumberField},
\lt{Projectivization}, \lt{EllipticCurve} --- developing heights over
fields with a theory of absolute values, and a
\lt{Mathlib/GroupTheory/Descent.lean} of M.~Stoll giving an abstract
descent theorem (a height function with suitable growth relative to an
endomorphism forces finite generation), one of the ingredients of
Mordell--Weil. The elliptic-curve height file states as its goal the
approximate parallelogram law $\lvert h(P{+}Q) + h(P{-}Q) - 2(h(P) +
h(Q))\rvert \le C$ for the na\"ive height, lists it as a
\textsc{todo}, and currently contains one theorem: that the
sum-and-difference map on $\mathbb{P}^2$ induced by a Weierstrass curve
is a morphism, in the form
$\exists C\, \forall x,\ \lvert \log\nh(F(x)) - 2\log\nh(x)\rvert \le
C$. That is the same estimate our \S\ref{sec:secant} proves, in far
greater generality of ground field --- and with a non-explicit
constant. The difference matters for our purposes: an unquantified $C$
cannot be fed to interval arithmetic, whereas the constants below are
polynomials in the coefficients and can be evaluated. Neither the
canonical height nor any rank bound is present there. Our development
is independent of both files.

Buzzard's FLT project \cite{FLT} targets a modularity lifting theorem
sufficient for Fermat, not the full modularity theorem. Gross--Zagier,
Kolyvagin and Heegner point theory are, to our knowledge, unformalized
anywhere.

\section{Setup and the \texorpdfstring{$b$}{b}-invariant basis}

Throughout, $W$ is a Weierstrass curve over $\mathbb{Q}$ with
coefficients $(a_1,a_2,a_3,a_4,a_6)$, and $W(\mathbb{Q})$ is
\texttt{mathlib}'s \lt{WeierstrassCurve.Affine.Point}: the inductive
type with constructors \lt{zero} and \lt{some}, the latter carrying a
nonsingular affine rational point, with \texttt{mathlib}'s
\lt{AddCommGroup} instance. The group law requires only that the ground
field have decidable equality; nonsingularity of the curve is not an
instance hypothesis, which is why $\Delta \ne 0$ appears explicitly in
Theorem~\ref{thm:main}.

The $b$-invariants are
\[
  b_2 = a_1^2 + 4a_2, \quad b_4 = 2a_4 + a_1a_3, \quad
  b_6 = a_3^2 + 4a_6, \quad
  b_8 = a_1^2a_6 + 4a_2a_6 - a_1a_3a_4 + a_2a_3^2 - a_4^2,
\]
subject to \texttt{mathlib}'s \lt{b\_relation}
\begin{equation}\label{eq:brel}
  4b_8 = b_2b_6 - b_4^2 ,
\end{equation}
and $\Delta = -b_2^2b_8 + 9b_2b_4b_6 - 8b_4^3 - 27b_6^2$. Hypothesis
(i) of Theorem~\ref{thm:main} says that the four $b$-invariants are the
images of integers $B_2,B_4,B_6,B_8$; the integral form of
\eqref{eq:brel} is then \emph{derived} from \texttt{mathlib}'s rational
one rather than assumed
(\lt{b\_relation\_int}). We drop the capitals below and write
$b_i$ for the integers.

Choosing the $b$-invariants as the ground basis, rather than the
$a_i$, is not cosmetic. Every object in this paper --- $\Phi$, $\Psi$,
$\Delta$, and the three secant forms --- is a polynomial in the $b_i$
alone. Expressed over $\mathbb{Z}[a_1,\dots,a_6]$ the B\'ezout
cofactors of \S\ref{sec:dup} would be unmanageable; over
$\mathbb{Z}[b_2,b_4,b_6,b_8]$ the four of them comprise $168$ monomials
in total ($55$, $54$, $26$ and $33$ respectively), of total degree at
most nine in $(X,Y,b_2,b_4,b_6,b_8)$, which the kernel checks by
\lt{ring} without complaint. In that free polynomial ring
$\Res(\Phi,\Psi)$ is not literally $\Delta^2$; the two agree modulo
\eqref{eq:brel}, and the correction multiplier is eight monomials, so
each identity closes with a single \lt{linear\_combination}.

\begin{definition}[\lt{naiveHeight}, \lt{nh}, \lt{X}]
For $q \in \mathbb{Q}$ set
$\nh(q) := \max\bigl(\lvert\mathrm{num}(q)\rvert, \mathrm{den}(q)\bigr)$,
using Lean's reduced representation; $\nh(q) \ge 1$ always. Define
$x \colon W(\mathbb{Q}) \to \mathbb{Q}$ by $x(O) = 0$ and
$x(\text{some at }(X,Y)) = X$, and $h(R) := \log \nh(x(R))$.
\end{definition}

\begin{remark}[Notation]\label{rem:notation}
In the Lean sources the two points' reduced coordinates are named
$(a_1,b_1,a_2,b_2)$, which collides typographically with the curve's
$a$- and $b$-invariants. Here we write $x_i = u_i/v_i$ with
$\gcd(u_i,v_i)=1$ and reserve $a_i, b_i$ for the invariants. The
translation is $u_i \leftrightarrow a_i$, $v_i \leftrightarrow b_i$ in
the arguments of \lt{DDz}, \lt{Sfz}, \lt{Prz} and their certificates.
\end{remark}

\section{The duplication side}\label{sec:dup}

This section records what \cite{MW} already had in universal form; it is
included because the secant side reduces to it.

\begin{proposition}[\lt{bezout\_X}, \lt{bezout\_Y},
\lt{dvd\_disc\_sq\_of\_isCoprime}]\label{prop:bez}
With $\Phi,\Psi$ as in \eqref{eq:PhiPsi} there are explicit
$F_1,G_1,F_2,G_2 \in \mathbb{Z}[b_2,b_4,b_6,b_8][X,Y]$, binary cubics in
$(X,Y)$, such that, granted \eqref{eq:brel},
\[
  F_1\Phi + G_1\Psi = \Delta^2 X^7, \qquad
  F_2\Phi + G_2\Psi = \Delta^2 Y^7 .
\]
Consequently, if $\gcd(X,Y) = 1$, every common divisor of $\Phi(X,Y)$
and $\Psi(X,Y)$ divides $\Delta^2$.
\end{proposition}

\begin{proposition}[\lt{max\_abs\_le}, \lt{disc\_sq\_mul\_le}]
\label{prop:dupsize}
With $M = \max(\lvert X\rvert,\lvert Y\rvert)$,
\[
  \max\bigl(\lvert\Phi\rvert,\lvert\Psi\rvert\bigr) \le CU\cdot M^4,
  \qquad
  \Delta^2 M^4 \le CL \cdot
    \max\bigl(\lvert\Phi\rvert,\lvert\Psi\rvert\bigr),
\]
where
\[
  CU = \max\bigl(1 + \lvert b_4\rvert + 2\lvert b_6\rvert +
        \lvert b_8\rvert,\;
        4 + \lvert b_2\rvert + 2\lvert b_4\rvert + \lvert b_6\rvert\bigr)
\]
is the coefficient sum and $CL$ is the corresponding sum of absolute
coefficients of the B\'ezout cofactors,
$CL = \max\bigl(S_{F_1}+S_{G_1},\, S_{F_2}+S_{G_2}\bigr)$.
\end{proposition}

The lower bound is Proposition~\ref{prop:bez} read archimedeanly: each
cofactor is a binary cubic, so
$\Delta^2M^7 \le CL\cdot M^3 \max(\lvert\Phi\rvert,\lvert\Psi\rvert)$,
and $M^3$ cancels. The same identity that removes the content analysis
removes the archimedean one.

\begin{proposition}[\lt{nh\_dblQ\_le},
\lt{nh\_pow\_le\_of\_disc\_ne\_zero}, \lt{addX\_self\_eq\_dblQ}]
\label{prop:dupheight}
Let $\mathrm{dbl}(x) := \Phi(u,v)/\Psi(u,v)$ for $x = u/v$ reduced. Then
$\mathrm{dbl}$ agrees with \texttt{mathlib}'s group law: for an affine
point $(x,y)$ with $2\cdot(x,y) \ne O$ one has
$x\bigl(2\cdot(x,y)\bigr) = \mathrm{dbl}(x)$. Moreover, if
$\Delta \ne 0$ and $\Psi(u,v) \ne 0$,
\[
  \nh\bigl(\mathrm{dbl}(x)\bigr) \le CU\cdot \nh(x)^4,
  \qquad
  \nh(x)^4 \le CL\cdot \nh\bigl(\mathrm{dbl}(x)\bigr).
\]
\end{proposition}

The unreduced fraction $\Phi/\Psi$ has content
$g = \gcd(\Phi,\Psi)$ and $g\cdot \nh(\Phi/\Psi) =
\max(\lvert\Phi\rvert,\lvert\Psi\rvert)$; then $g \ge 1$ gives the upper
bound and $g \le \Delta^2$ (Proposition~\ref{prop:bez}) gives the lower
one, the $\Delta^2$ cancelling against the $\Delta^2$ of
Proposition~\ref{prop:dupsize}. The duplication formula itself is
Silverman \cite[III.2.3(d)]{silverman}, and is proved here against
\texttt{mathlib}'s \lt{addX}/\lt{slope} --- a polynomial identity modulo
the curve equation whose certificate is the single multiplier
$-(b_2 + 8x)$ --- rather than assumed.

Setting $\kappa := \max(CU, CL)$, Propositions~\ref{prop:dupsize}
and~\ref{prop:dupheight} give the logarithmic doubling window
$\lvert h(2R) - 4h(R)\rvert \le \log\kappa$ away from $O$, and at $O$
both sides vanish by the convention $x(O) = 0$
(\lt{heightWindow\_of\_xdbl\_ne\_zero}). That is the sole input of the
generic construction of \cite{MW}: given an abelian group $G$, a
nonnegative $\lambda \colon G \to \mathbb{R}$ and $\kappa \ge 1$ with
$\lvert \lambda(R+R) - 4\lambda(R)\rvert \le \log\kappa$, the sequence
$\lambda(2^nR)/4^n$ is Cauchy with geometric modulus and its limit
$\hh$ satisfies $\hh \ge 0$, $\hh(R+R) = 4\hh(R)$,
$\hh(2^nR) = 4^n\hh(R)$, $\lvert\hh - \lambda\rvert \le \log\kappa/3$
and $\lvert\hh(R) - \lambda(2^nR)/4^n\rvert \le \log\kappa/(3\cdot4^n)$.

\section{The secant side}\label{sec:secant}

\subsection{The universal formulas}

Both $x(P{+}Q) + x(P{-}Q)$ and $x(P{+}Q)\cdot x(P{-}Q)$ are symmetric in
the $y$-coordinates, so both eliminate them.

\begin{proposition}[\lt{addX\_add\_addX\_negY},
\lt{addX\_mul\_addX\_negY}]\label{prop:secform}
Over any field, for points $P = (x_1,y_1)$, $Q = (x_2,y_2)$ on a
Weierstrass curve with $x_1 \ne x_2$,
\[
  x(P{+}Q) + x(P{-}Q)
    = \frac{2x_1x_2(x_1+x_2) + b_2x_1x_2 + b_4(x_1+x_2) + b_6}
           {(x_1-x_2)^2},
\]
\[
  x(P{+}Q)\cdot x(P{-}Q)
    = \frac{x_1^2x_2^2 - b_4x_1x_2 - b_6(x_1+x_2) - b_8}
           {(x_1-x_2)^2}.
\]
\end{proposition}

Both are proved against \texttt{mathlib}'s \lt{addX} and \lt{slope} by
\lt{linear\_combination} modulo the two curve equations. The
$b$-invariant basis is again what makes the certificates small: the sum
identity has cofactors $2,2$.

\subsection{The integer forms}

Homogenizing Proposition~\ref{prop:secform} at $x_i = u_i/v_i$ gives
three integer forms of bidegree $(2,2)$:
\begin{equation}\label{eq:secforms}
\begin{aligned}
  \DDz &= (u_1v_2 - u_2v_1)^2,\\
  \Sfz &= 2u_1u_2(u_1v_2 + u_2v_1) + b_2u_1u_2v_1v_2
          + b_4(u_1v_2+u_2v_1)v_1v_2 + b_6v_1^2v_2^2,\\
  \Prz &= u_1^2u_2^2 - b_4u_1u_2v_1v_2
          - b_6(u_1v_2+u_2v_1)v_1v_2 - b_8v_1^2v_2^2 .
\end{aligned}
\end{equation}

\begin{proposition}[\lt{DDz\_spec}, \lt{Sfz\_spec}, \lt{Prz\_spec},
\lt{sum\_relation}, \lt{prod\_relation}]\label{prop:spec}
The bihomogenization is exact: with $x_3 = x(P{+}Q)$ and
$x_4 = x(P{-}Q)$,
\[
  \DDz\,(x_3 + x_4) = \Sfz, \qquad \DDz\, x_3x_4 = \Prz ,
\]
as identities in $\mathbb{Q}$. In particular $x_3$ and $x_4$ are the two
roots of the integer quadratic
$\DDz\,T^2 - \Sfz\,T + \Prz$.
\end{proposition}

\subsection{Diagonal degeneration}\label{sec:diag}

\begin{proposition}[\lt{Sfz\_diag}, \lt{Prz\_diag}, \lt{DDz\_diag}]
\label{prop:diag}
$\Sfz(u,v,u,v) = \Psi(u,v)$, $\Prz(u,v,u,v) = \Phi(u,v)$, and
$\DDz(u,v,u,v) = 0$.
\end{proposition}

\begin{proof}[Proof, as formalized]
Substitute and expand:
$2u^2(2uv) + b_2u^2v^2 + b_4(2uv)v^2 + b_6v^4
 = 4u^3v + b_2u^2v^2 + 2b_4uv^3 + b_6v^4 = \Psi(u,v)$, and
$u^4 - b_4u^2v^2 - b_6(2uv)v^2 - b_8v^4 = \Phi(u,v)$. Each is one
\lt{ring} call.
\end{proof}

Proposition~\ref{prop:diag} is the structural observation of this paper.
It says that the secant construction, evaluated where the secant line
degenerates to the tangent, returns the duplication construction --- not
up to a unit, not up to the relation \eqref{eq:brel}, but on the nose.
It is what one would hope for, and it is worth stating because of its
consequence: the content problem for the secant triple can be
\emph{reduced} to the content problem for the duplication pair, and the
latter has already been solved once and for all by
$\Res(\Phi,\Psi) = \Delta^2$. There is no second resultant to compute.
The reduction is carried out in the next subsection; the certificates
that carry it out are exactly the ones that produce
$\Phi^2$ and $\Psi^2$ from eliminating the second point, which is the
$x_2$-resultant statement $\Res_{x_2}(\DDz,\Sfz) = \psi(x_1)^2$,
$\Res_{x_2}(\DDz,\Prz) = \varphi(x_1)^2$ in bihomogenized form.

\subsection{The content bound}\label{sec:content}

\begin{proposition}[\lt{psi\_sq\_cert}, \lt{phi\_sq\_cert},
\lt{psi\_sq\_certA}, \lt{phi\_sq\_certA}]\label{prop:twin}
There are explicit integer cofactors
$\sigma_1,\sigma_2,\sigma_1',\sigma_2'$ and
$\pi_1,\pi_2,\pi_1',\pi_2'$ in
$\mathbb{Z}[b_2,b_4,b_6,b_8][u_1,v_1,u_2,v_2]$ with
\[
  \sigma_1\DDz + \sigma_2\Sfz = v_2^3\,\Psi(u_1,v_1)^2, \qquad
  \sigma_1'\DDz + \sigma_2'\Sfz = u_2^3\,\Psi(u_1,v_1)^2,
\]
\[
  \pi_1\DDz + \pi_2\Prz = v_2^3\,\Phi(u_1,v_1)^2, \qquad
  \pi_1'\DDz + \pi_2'\Prz = u_2^3\,\Phi(u_1,v_1)^2 .
\]
All eight are verified by \lt{ring}.
\end{proposition}

The pairing of weights is the point. Sylvester's matrix, cleared of
denominators, delivers the resultant multiplied by a power of the
leading coefficient, and which coefficient is leading depends on which
variable one eliminates; carrying both eliminations gives the same
target with weights $u_2^3$ and $v_2^3$. Since $\gcd(u_2,v_2) = 1$ we
have $\gcd(u_2^3,v_2^3) = 1$, and a B\'ezout combination cancels the
weights (\lt{dvd\_of\_dvd\_coprime\_weights}). Either certificate
without its twin leaves an uncontrolled monomial factor and is useless.

\begin{theorem}[\lt{secant\_content\_universal}]\label{thm:content}
Let $b_2,b_4,b_6,b_8 \in \mathbb{Z}$ satisfy \eqref{eq:brel}, and let
$u_1/v_1$, $u_2/v_2$ be in lowest terms. Then every common divisor of
$\DDz$, $\Sfz$, $\Prz$ divides $\Delta^6$.
\end{theorem}

\begin{proof}[Proof, as formalized]
Let $d$ divide all three. By Proposition~\ref{prop:twin} and weight
cancellation, $d \mid \Psi(u_1,v_1)^2$ and $d \mid \Phi(u_1,v_1)^2$.
Apply \eqref{eq:cube} to the B\'ezout identity
$F_1\Phi + G_1\Psi = \Delta^2u_1^7$ of Proposition~\ref{prop:bez}: every
term of the cube lies in the ideal generated by $\Phi^2$ and $\Psi^2$,
so $d \mid (\Delta^2u_1^7)^3 = u_1^{21}\Delta^6$; the $Y$-side identity
gives $d \mid v_1^{21}\Delta^6$; and $\gcd(u_1,v_1)=1$ cancels the
monomials.
\end{proof}

This is the uniform replacement for what \cite{MW} did per curve: for
\texttt{389a1} an extended Euclidean computation, and for
\texttt{5077a1} the same plus a Hermite-normal-form argument to show
that a stray factor was forced. Both are subsumed. The exponent $6$ is
the price of the cube in \eqref{eq:cube}; sharpening it for a particular
curve remains possible and is not needed downstream, because the content
bound is only ever used to know that a content exists and is bounded.

\subsection{Size bounds}\label{sec:size}

\begin{proposition}[\lt{abs\_DDz\_le}, \lt{abs\_Sfz\_le},
\lt{abs\_Prz\_le}]\label{prop:secupper}
With $H_i = \max(\lvert u_i\rvert, \lvert v_i\rvert)$,
\[
  \lvert\DDz\rvert \le 4H_1^2H_2^2, \quad
  \lvert\Sfz\rvert \le \bigl(4 + \lvert b_2\rvert + 2\lvert b_4\rvert
     + \lvert b_6\rvert\bigr)H_1^2H_2^2, \quad
  \lvert\Prz\rvert \le \bigl(1 + \lvert b_4\rvert + 2\lvert b_6\rvert
     + \lvert b_8\rvert\bigr)H_1^2H_2^2 .
\]
\end{proposition}

These are triangle inequalities on bidegree-$(2,2)$ monomials, and they
are sharp against the hand computations of \cite{MW}: at
\texttt{5077a1} the $\Sfz$ constant is $4+0+28+25 = 57$ and the $\Prz$
constant $1+14+50+49 = 114$, which are the numbers derived by hand
there. Writing
\begin{equation}\label{eq:CUsec}
  CU_{\mathrm{sec}} := \max\bigl(4,\;
    4 + \lvert b_2\rvert + 2\lvert b_4\rvert + \lvert b_6\rvert,\;
    1 + \lvert b_4\rvert + 2\lvert b_6\rvert + \lvert b_8\rvert\bigr)
   \;=\; \max(4, CU),
\end{equation}
all three are bounded by $CU_{\mathrm{sec}}H_1^2H_2^2$.

There is a matching lower bound, proved from the same certificates in
the manner of Proposition~\ref{prop:dupsize}.

\begin{proposition}[\lt{secant\_size\_lower}]\label{prop:seclower}
There is an explicit $CL_{\mathrm{sec}} = CL^2\cdot S$, with $S$ the sum
of the coefficient-absolute-sums of the eight cofactors of
Proposition~\ref{prop:twin}, such that
\[
  \Delta^4 H_1^2H_2^2 \;\le\; CL_{\mathrm{sec}}\cdot
    \max\bigl(\lvert\DDz\rvert, \lvert\Sfz\rvert,
              \lvert\Prz\rvert\bigr).
\]
\end{proposition}

\begin{remark}
Proposition~\ref{prop:seclower} is the honest analogue of
Proposition~\ref{prop:dupsize} on the secant side --- two-sided control
of the forms themselves --- but it is \emph{not} used in the
parallelogram assembly below. The lower half of the
quasi-parallelogram law comes for free from the doubling trick
(\S\ref{sec:window}), and the plan that produced
Proposition~\ref{prop:seclower} had assumed otherwise. We record the
correction rather than quietly presenting the result as load-bearing.
\end{remark}

\section{The uniform height window}\label{sec:window}

\begin{lemma}[\lt{naiveHeight\_mul\_le\_of\_quadratic}]\label{lem:roots}
Let $A,B,C \in \mathbb{Z}$ with $A \ne 0$ and let $u,v \in \mathbb{Q}$
satisfy $A(u+v) = B$ and $Auv = C$. Then
$\nh(u)\nh(v) \le 2\max(\lvert A\rvert,\lvert B\rvert,\lvert C\rvert)$.
\end{lemma}

Lemma~\ref{lem:roots} is curve-independent and was already available
from \cite{MW}, where it is proved without Gauss's lemma. Combining it
with Propositions~\ref{prop:spec} and~\ref{prop:secupper} gives the
upper half of the window at the point level.

\begin{proposition}[\lt{secant\_upper\_window}]\label{prop:upper}
On any Weierstrass curve over $\mathbb{Q}$ with integer $b$-invariants,
for affine points with $x_1 \ne x_2$,
\[
  \nh\bigl(x(P{+}Q)\bigr)\,\nh\bigl(x(P{-}Q)\bigr)
    \;\le\; 2\,CU_{\mathrm{sec}}\;\nh(x_1)^2\nh(x_2)^2 .
\]
\end{proposition}

The lower half is not a second estimate but the same one, applied
elsewhere. The sum and difference of the pair $\bigl(P{+}Q, P{-}Q\bigr)$
are $2P$ and $2Q$; so Proposition~\ref{prop:upper} at that pair, chained
with the duplication lower window of
Proposition~\ref{prop:dupheight}, bounds
$\nh(x_1)^2\nh(x_2)^2$ from above. Doing so requires knowing
$x(P{+}Q) \ne x(P{-}Q)$, which is where hypothesis (iii) first bites:
equality would force $P{+}Q = \pm(P{-}Q)$, hence $2Q = O$ or $2P = O$
(\lt{X\_add\_ne\_X\_sub}).

\begin{theorem}[\lt{point\_two\_sided}, \lt{lognhX\_defect\_le}]
\label{thm:twosided}
Let $W$ have integer $b$-invariants and $\Delta \ne 0$, and let $P,Q$ be
affine points with $x_1 \ne x_2$, neither $2$-torsion. Put
\begin{equation}\label{eq:Cpar}
  C_{\mathrm{par}} := \bigl(CL^2 + 1\bigr)\cdot
     2\,CU_{\mathrm{sec}} .
\end{equation}
Then
\[
  \nh(x_1)^2\nh(x_2)^2
    \;\le\; C_{\mathrm{par}}\,
      \nh\bigl(x(P{+}Q)\bigr)\nh\bigl(x(P{-}Q)\bigr)
    \;\le\; C_{\mathrm{par}}\cdot 2\,CU_{\mathrm{sec}}\,
      \nh(x_1)^2\nh(x_2)^2,
\]
equivalently
$\bigl\lvert h(P{+}Q) + h(P{-}Q) - 2h(P) - 2h(Q)\bigr\rvert
 \le \log C_{\mathrm{par}}$.
\end{theorem}

The $+1$ in \eqref{eq:Cpar} is there so that
$C_{\mathrm{par}} \ge 2CU_{\mathrm{sec}}$ unconditionally and one
constant serves both sides of the window; it costs nothing and removes a
case split.

\begin{theorem}[\lt{heightWindow\_universal}]\label{thm:windowuniv}
Let $W$ be admissible. Then $R \mapsto \log\nh(x(R))$ satisfies the
doubling window with constant $\kappa = \max(CU,CL)$ on all of
$W(\mathbb{Q})$, with no per-curve input: the integral $b$-relation is
derived from \texttt{mathlib}'s, the nonvanishing of $\Psi$ comes from
(iii), and the identification of $x(2R)$ with $\Phi/\Psi$ comes from
\texttt{mathlib}'s group law.
\end{theorem}

Theorem~\ref{thm:windowuniv} is what makes $\hh$ exist on an admissible
curve, by the generic construction recalled at the end of
\S\ref{sec:dup}. We write $\hh = $ \lt{canheightW} $W$.

\section{The parallelogram law and bi-additivity}\label{sec:par}

\begin{theorem}[\lt{canheight\_parallelogram\_universal}]\label{thm:parx}
Let $W$ be admissible and let $P,Q$ be affine points with distinct
$x$-coordinates. Then
$\hh(P{+}Q) + \hh(P{-}Q) = 2\hh(P) + 2\hh(Q)$.
\end{theorem}

\begin{proof}[Proof, as formalized]
Hypothesis (iii) forces $2^nR \ne O$ for every $n$ and every $R \ne O$,
and forces $x(2^nP) \ne x(2^nQ)$ at every level (else $2^n(P \mp Q) =
O$). Theorem~\ref{thm:twosided} at $(2^nP, 2^nQ)$, divided by $4^n$,
bounds the defect of the scaled sequences by $\log
C_{\mathrm{par}}/4^n \to 0$; the scaled sequences converge to the four
canonical heights; uniqueness of limits closes it.
\end{proof}

\begin{theorem}[\lt{par}]\label{thm:parall}
On an admissible curve the parallelogram law holds for \emph{every} pair
$P,Q \in W(\mathbb{Q})$, with no side conditions.
\end{theorem}

\begin{proof}[Proof, as formalized]
The degenerate cases $P = O$, $Q = O$, $Q = -P$, $P = Q$ are closed by
$\hh(O) = 0$, $\hh(-R) = \hh(R)$ and $\hh(2R) = 4\hh(R)$. Otherwise both
points are affine, and equal $x$-coordinates would force one of the
degenerate cases, since two affine points of a Weierstrass curve with
the same $x$-coordinate are equal or negatives
(\lt{eq\_or\_eq\_neg\_of\_x\_eq}). So distinct $x$-coordinates may be
assumed and Theorem~\ref{thm:parx} applies. Here $\hh(-R) = \hh(R)$ is
hypothesis-free: negation fixes the $x$-coordinate, so the two defining
sequences are literally equal.
\end{proof}

\begin{corollary}[\lt{mul}]\label{cor:mult}
$\hh(kR) = k^2\hh(R)$ for every $k \in \mathbb{Z}$ and every
$R \in W(\mathbb{Q})$.
\end{corollary}

\begin{proof}[Proof, as formalized]
Theorem~\ref{thm:parall} at $\bigl((n{+}2)R, R\bigr)$ gives the
three-term recurrence
$\hh((n{+}3)R) = 2\hh((n{+}2)R) + 2\hh(R) - \hh((n{+}1)R)$; with base
values $\hh(R)$ and $\hh(2R) = 4\hh(R)$ a two-step induction gives all
$k \in \mathbb{N}$, and $\hh(-R) = \hh(R)$ gives the rest. Because
Theorem~\ref{thm:parall} is unconditional the recurrence needs no
indexing shift, unlike its conditional counterpart in \cite{MW}, where
the step at $k=0$ had a torsion difference point and had to be supplied
separately.
\end{proof}

\begin{theorem}[\lt{pairing\_add\_left}, \lt{pairing\_comm},
\lt{biAdditive\_pairingW}]\label{thm:biadd}
On an admissible curve the pairing
$\langle P,Q\rangle = \tfrac12\bigl(\hh(P{+}Q) - \hh(P) -
\hh(Q)\bigr)$ is symmetric and additive in each variable, and
$\langle P,P\rangle = \hh(P)$.
\end{theorem}

The derivation is Jordan--von Neumann: from
$4\langle P,Q\rangle = \hh(P{+}Q) - \hh(P{-}Q)$ one takes the four
instances of Theorem~\ref{thm:parall} at $(P{\pm}R, Q)$ and
$(Q{\pm}R, P)$, pairs the mixed differences using $\hh(-u) = \hh(u)$,
and adds. Symmetry is immediate from the definition. Because
Theorem~\ref{thm:parall} is unconditional, no side condition has to be
propagated through the four instances --- which is precisely the
bookkeeping that \cite{MW} avoided, for two points, by a different route
(a hypothesized relation reduces everything to multiples of one point),
and paid for, for three points, by proving torsion-freeness.

\section{The rank pipeline}\label{sec:pipeline}

The last step is pure linear algebra and was already curve-independent
in \cite{MW}: if $B$ is a symmetric bi-additive real form on an abelian
group $G$ and $P_1,\dots,P_n \in G$ have
$\det\bigl(B(P_i,P_j)\bigr) \ne 0$, then
$n \le \rank_{\mathbb{Z}} G$ (\lt{rank\_ge\_of\_gram\_det\_ne\_zero}). A
relation $\sum v_iP_i = 0$ pairs against each $P_j$ to put $v$ in the
kernel of the Gram matrix; a matrix over a field with a nonzero kernel
vector has zero determinant; so the map
$\mathbb{Z}^n \to G$, $v \mapsto \sum v_iP_i$, is injective.

Theorem~\ref{thm:main} is Theorem~\ref{thm:biadd} fed into that
statement, and its proof term is two lines. The two specializations
\eqref{eq:reg2} and \lt{gram\_det\_fin\_three} are the $2\times2$ and
$3\times3$ expansions, obtained by \lt{Matrix.det\_fin\_two} /
\lt{Matrix.det\_fin\_three} together with
$\langle P,P\rangle = \hh(P)$ and symmetry.

It is worth being explicit about what the pipeline consumes and what it
does not produce. It consumes: four integers, a proof that they are the
$b$-invariants of the curve, a proof that $\Delta \ne 0$, a proof of no
rational $2$-torsion, a list of points, and a proof that a determinant
is nonzero. It does not produce the last of these. Certifying
$\det \ne 0$ for a given curve is an interval-arithmetic computation,
done per curve; see \S\ref{sec:notproved}(c).

\section{Validation, and a definitional collapse}\label{sec:validation}

\subsection{The two curves}

For $E = \mathtt{389a1}$ the hypothesis package is discharged as
follows: the $b$-invariants $(4,-4,1,-3)$ by \lt{norm\_num} from
$(a_1,\dots,a_6) = (0,1,1,-2,0)$; $\Delta = 389 \ne 0$ by \lt{decide} on
the integer polynomial; and (iii) from the triviality of
$E(\mathbb{Q})_{\mathrm{tors}}$, proved in \cite{MW}. For
$E' = \mathtt{5077a1}$, with $(a_1,\dots,a_6) = (0,0,1,-7,6)$, the
$b$-invariants are $(0,-14,25,-49)$ and $\Delta = 5077$, again by
\lt{decide}, and (iii) again from triviality of the torsion subgroup.
Feeding in the certified regulator determinants of \cite{MW} yields
\eqref{eq:389} and \eqref{eq:5077}. For the three-point case one needs
that the $3\times3$ Gram determinant of the uniform pairing is the
determinant computed in \cite{MW}; that is
\lt{gram\_det\_eq\_regDet3}, closed by \lt{ring} after rewriting, with
no slack.

\subsection{The definitional collapse}

Transporting the numerical certificates required a bridge, because the
two layers reach the canonical height through different helper
functions: the uniform layer through
$\log\nh$ composed with a generic $x$-projection, the per-curve layer of
\cite{MW} through $\log$ of a differently-defined na\"ive height
composed with a per-curve $x$-projection. The two $x$-projections agree
by \lt{cases} then \lt{rfl} (the pattern matches are the same), and the
two na\"ive heights agree by a cast. This gives the two log-height
helpers as \emph{equal functions}.

The point worth recording is what happens next. Once the helpers are
equal as functions, the two approximating sequences and hence the two
\lt{limUnder}s are \emph{definitionally} equal:
\lt{canheightW\_eq\_389} closes by \lt{rfl}, and so does
\lt{canheightW\_eq\_5077}. No analytic argument about limits was needed
and none was attempted.

\begin{remark}\label{rem:collapse}
We think this is a real point about formalized refactoring rather than
an anecdote. The risk in a universalization of this depth is not that a
proof breaks --- broken proofs announce themselves --- but that the
object being computed silently changes, so that the new theorems are
true of something slightly different from the old ones and the
comparison is never actually made. Here the comparison is made, and it
is made by the kernel, at the strongest available level: not
propositional equality proved by an argument, but definitional equality
checked by reduction. The uniform layer and the per-curve layer are the
same construction.
\end{remark}

\section{The price of uniformity}\label{sec:price}

The universal constants are much worse than the hand-built ones. The
measured values, for the two curves of \cite{MW}, are the following.

\begin{center}
\small
\begin{tabular}{lcrrrrrr}
\toprule
curve & $(b_2,b_4,b_6,b_8)$ & $\Delta$ & $CU$ & $CL = \kappa$
  & $\kappa$ from \cite{MW} & $CU_{\mathrm{sec}}$
  & $C_{\mathrm{par}}$ \\
\midrule
\texttt{389a1}  & $(4,-4,1,-3)$    & $389$  & $17$  & $672\,192$
  & $1\,728$   & $17$  & ${\approx}1.54\times10^{13}$\\
\texttt{5077a1} & $(0,-14,25,-49)$ & $5077$ & $114$ & $536\,913\,058$
  & $105\,754$ & $114$ & ${\approx}6.57\times10^{19}$\\
\bottomrule
\end{tabular}
\end{center}

\noindent
Exactly, $C_{\mathrm{par}} = 15\,362\,630\,885\,410$ for
\texttt{389a1} and
$C_{\mathrm{par}} = 65\,726\,844\,062\,007\,791\,220$ for
\texttt{5077a1}. Every entry is obtained by evaluating the definitions
of \S\S\ref{sec:dup}--\ref{sec:window} at the four integers in the
second column; the $\kappa$ of \cite{MW} is that paper's per-curve
duplication constant.

Three remarks, in decreasing order of comfort.

\emph{The upper side is sharp.} $CU = 17$ and $CU = 114$ are exactly the
size constants derived by hand in \cite{MW} for the two curves. Nothing
is lost on the $CU$ side, and hence nothing on $CU_{\mathrm{sec}}$,
which by \eqref{eq:CUsec} is $\max(4,CU)$.

\emph{The loss is entirely in $CL$, and it is one factor of $\Delta$.}
The measured ratios are
$672\,192/1\,728 = 389$ and $536\,913\,058/105\,754 = 5077$ --- in both
cases exactly $\Delta$. Neither constant was designed to make this
happen. The natural reading is that the universal route normalizes its
B\'ezout identity to the resultant, $F\Phi + G\Psi = \Delta^2X^7$, and
then cancels a content bounded by $\Delta^2$; whereas for these two
curves the actual content of $(\Phi,\Psi)$ is only $\Delta$, and the
hand computation of \cite{MW} found the sharper identity with $\Delta$
on the right. The uniform statement cannot know that, because
$\Res(\Phi,\Psi) = \Delta^2$ is what is true in the free polynomial
ring. Beyond that structural factor, $CL$ is a sum of absolute values of
B\'ezout cofactors, which is not tight.

\emph{It is still the right trade.} The cost is visible downstream: the
torsion-enumeration argument of \cite{MW} searches
$\kappa^{1/3}$-many candidates, so for \texttt{389a1} the universal
constant would widen the search from $12$ to about $88$; and
$\log C_{\mathrm{par}}$ enters the parallelogram defect, which is
squeezed by $4^{-n}$, so matching the accuracy that \cite{MW}'s
constant $10368$ gives at \texttt{389a1} would need about
$\log_4\bigl(C_{\mathrm{par}}/10368\bigr) \approx 15$ further dyadic
levels. Both are real costs, and both are
costs in \emph{numerics}, not in \emph{proof}. Against them: a constant
that is automatic for every curve, obtained by evaluating a polynomial
at four integers, versus a sharp constant that costs a hand-built
certificate suite per curve. Sharpening one curve later is optional;
deriving a constant at all is not. Nothing above prevents a user from
supplying a better per-curve $\kappa$ to the same generic machinery ---
that is exactly how \cite{MW}'s numerics were produced --- and the
uniform bound remains as the fallback that always exists.

\section{Outlook}

Three things now cost less than they did.

\emph{A new curve.} The scale-out cost is the hypothesis package ---
four $b$-invariants, $\Delta \ne 0$, no rational $2$-torsion --- plus a
certified-nonzero Gram determinant at the chosen points. The first three
are \lt{decide}-level once the coefficients are fixed; the last is
interval arithmetic. Nothing else.

\emph{Higher rank.} Theorem~\ref{thm:main} is stated for all $n$, so
rank $\ge 4$ is a matter of supplying four points and one determinant.
The algebra is finished.

\emph{Upstreaming.} Three pieces are of general interest and are
independent of everything else here: Lemma~\ref{lem:roots}; the
implication ``two-sided duplication estimate $\Rightarrow$ canonical
height''; and the Gram criterion. A fourth, more speculative candidate
is the pair (Proposition~\ref{prop:diag},
Proposition~\ref{prop:twin}) --- the diagonal degeneration and the twin
certificates --- which is what turns the approximate parallelogram law
from a per-curve computation into a theorem with explicit constants.
\texttt{mathlib}'s \lt{NumberTheory/Height/EllipticCurve.lean} lists the
approximate parallelogram law as an open task
(\S\ref{sec:related}); the route taken here is one way to it, at the
cost of specializing the ground field to $\mathbb{Q}$ and buying
explicit constants in exchange.

\section{Verification data}\label{sec:verif}

All results were checked with Lean~4 toolchain
\texttt{leanprover/lean4:v4.24.0-rc1} against \texttt{mathlib} at commit
\ltw{eed770a434957369c6262aa3fb1d6426419016d4}. At repository commit
\ltw{b99f2ab6abee7f6df627a4f070122b71882dac52} (5~August 2026),
\lt{lake build PF} completes in $4711$ jobs with zero errors and zero
warnings.

The one-command reproduction is
\begin{center}
\lt{cd PF\_Lean4\_Code \&\& lake build PF}
\end{center}
Each source file listed below ends with \lt{\#print axioms}
declarations for its own results, so the axiom report is printed by that
same build and requires no separate step. Every theorem cited in this
paper reports a subset of
\[
  [\mathtt{propext},\ \mathtt{Classical.choice},\ \mathtt{Quot.sound}]
\]
--- several of the polynomial identities report only
$[\mathtt{propext}]$ or $[\mathtt{propext}, \mathtt{Quot.sound}]$. The
development contains no \lt{sorry}, no \lt{native\_decide} (which would
introduce \lt{Lean.ofReduceBool}), and no project-level \lt{axiom}
declarations. We do not report a wall-clock figure: the build was
replayed from cache for this measurement, so any time we could quote
would not be the from-scratch time a reader would see.

The table below maps each claim to its Lean declaration. All files are
in \lt{PF\_Lean4\_Code/PF/}; file names carry an internal revision tag,
which is not part of the mathematics.

\begingroup
\footnotesize
\setlength{\LTleft}{0pt}\setlength{\LTright}{0pt}
\setlength{\tabcolsep}{4pt}
\begin{longtable}{@{}p{0.24\textwidth}p{0.34\textwidth}p{0.38\textwidth}@{}}
\toprule
claim & Lean declaration & file \\
\midrule
\endfirsthead
\toprule
claim & Lean declaration & file \\
\midrule
\endhead
\bottomrule
\endfoot
Prop.~\ref{prop:bez} (B\'ezout)
  & \lt{bezout\_X}, \lt{bezout\_Y}
  & \lt{DuplicationBezoutUniversal\_r174}\\
Prop.~\ref{prop:bez} (content)
  & \lt{dvd\_disc\_sq\_of\_isCoprime}
  & \lt{DuplicationBezoutUniversal\_r174}\\
Prop.~\ref{prop:dupsize}
  & \lt{max\_abs\_le}, \lt{disc\_sq\_mul\_le}
  & \lt{DuplicationSizeUniversal\_r175}\\
Prop.~\ref{prop:dupheight} (heights)
  & \lt{nh\_dblQ\_le},\newline
    \lt{nh\_pow\_le\_of\_disc\_ne\_zero}
  & \lt{DuplicationHeightUniversal\_r176}\\
Prop.~\ref{prop:dupheight} (formula)
  & \lt{addX\_self\_eq\_dblQ}, \lt{psiX\_eq\_sq}
  & \lt{DuplicationFormula\_r178}\\
doubling window
  & \lt{heightWindow\_of\_xdbl\_ne\_zero}
  & \lt{DuplicationLogWindowFixed\_r179}\\
generic canonical height
  & \lt{HeightWindow}, \lt{canheight},\newline
    \lt{canheight\_window\_shift}
  & \lt{CanonicalHeightGeneric\_r171}\\
Prop.~\ref{prop:secform}
  & \lt{addX\_add\_addX\_negY},\newline
    \lt{addX\_mul\_addX\_negY}
  & \lt{SecantUniversal\_r181}\\
\eqref{eq:secforms}, Prop.~\ref{prop:spec}
  & \lt{DDz}, \lt{Sfz}, \lt{Prz},\newline
    \lt{Sfz\_spec}, \lt{Prz\_spec}
  & \lt{SecantSizeUniversal\_r193}\\
Prop.~\ref{prop:secupper}
  & \lt{abs\_DDz\_le}, \lt{abs\_Sfz\_le},\newline
    \lt{abs\_Prz\_le}
  & \lt{SecantSizeUniversal\_r193}\\
Prop.~\ref{prop:diag}
  & \lt{Sfz\_diag}, \lt{Prz\_diag}, \lt{DDz\_diag}
  & \lt{SecantContentBridge\_r195}\\
Prop.~\ref{prop:twin} ($v_2^3$ twins)
  & \lt{psi\_sq\_cert}, \lt{phi\_sq\_cert}
  & \lt{SecantContentBridge\_r195}\\
Prop.~\ref{prop:twin} ($u_2^3$ twins)
  & \lt{psi\_sq\_certA}, \lt{phi\_sq\_certA}
  & \lt{SecantContentUniversal\_r196}\\
\eqref{eq:cube}, weight cancellation
  & \lt{dvd\_cube\_of\_bezout},\newline
    \lt{dvd\_of\_dvd\_coprime\_weights}
  & \lt{SecantContentUniversal\_r196}\\
Thm.~\ref{thm:content}
  & \lt{secant\_content\_universal}
  & \lt{SecantContentUniversal\_r196}\\
Prop.~\ref{prop:seclower}
  & \lt{secant\_size\_lower}, \lt{CLsec}
  & \lt{SecantSizeLower\_r197}\\
Lem.~\ref{lem:roots}
  & \lt{naiveHeight\_mul\_le\_of\_quadratic}
  & \lt{QuadraticRootHeights\_r148e}\\
Prop.~\ref{prop:spec} (integer form)
  & \lt{sum\_relation}, \lt{prod\_relation}
  & \lt{SecantUpperWindow\_r198}\\
Prop.~\ref{prop:upper}, \eqref{eq:CUsec}
  & \lt{secant\_upper\_window}, \lt{CUsec}
  & \lt{SecantUpperWindow\_r198}\\
Thm.~\ref{thm:twosided}, \eqref{eq:Cpar}
  & \lt{point\_two\_sided},\newline
    \lt{lognhX\_defect\_le}, \lt{Cpar}
  & \lt{PointParallelogramUniversal\_r199}\\
$x(P{+}Q) \ne x(P{-}Q)$
  & \lt{X\_add\_ne\_X\_sub},\newline
    \lt{eq\_or\_eq\_neg\_of\_x\_eq}
  & \lt{PointParallelogramUniversal\_r199}\\
Thm.~\ref{thm:windowuniv}, \eqref{eq:brel}
  & \lt{heightWindow\_universal},\newline
    \lt{b\_relation\_int}
  & \lt{CanheightParallelogramUniversal\_r200}\\
Thm.~\ref{thm:parx}
  & \lt{canheight\_parallelogram\_universal}
  & \lt{CanheightParallelogramUniversal\_r200}\\
Thm.~\ref{thm:height}(5), pairing
  & \lt{pairingW}, \lt{pairing\_self\_universal}
  & \lt{CanheightParallelogramUniversal\_r200}\\
Thm.~\ref{thm:parall}, Cor.~\ref{cor:mult}
  & \lt{par}, \lt{mul}, \lt{canheight\_neg},\newline
    \lt{canheight\_zero}
  & \lt{BiAdditiveUniversal\_r201}\\
Thm.~\ref{thm:biadd}
  & \lt{pairing\_add\_left}, \lt{pairing\_comm},\newline
    \lt{biAdditive\_pairingW}
  & \lt{BiAdditiveUniversal\_r201}\\
Gram criterion
  & \lt{rank\_ge\_of\_gram\_det\_ne\_zero}
  & \lt{GramRankGeneral\_r170}\\
Thm.~\ref{thm:main}
  & \lt{rank\_ge\_universal}
  & \lt{RankPipelineUniversal\_r202}\\
\eqref{eq:reg2} and $n=3$
  & \lt{gram\_det\_fin\_two},\newline
    \lt{rank\_ge\_two\_universal},\newline
    \lt{gram\_det\_fin\_three}
  & \lt{RankPipelineUniversal\_r202}\\
\S\ref{sec:validation}, bridges
  & \lt{canheightW\_eq\_389},\newline
    \lt{canheightW\_eq\_5077},\newline
    \lt{pairingW\_eq\_389}, \lt{gram\_det\_eq\_regDet3}
  & \lt{RankTransferWitnesses\_r203}\\
\eqref{eq:389}, \eqref{eq:5077}
  & \lt{E389a1\_rank\_ge\_two\_universal},\newline
    \lt{E5077a1\_rank\_ge\_three\_universal}
  & \lt{RankTransferWitnesses\_r203}\\
\end{longtable}
\endgroup

The B\'ezout cofactors of Proposition~\ref{prop:bez}, the size constants
of Proposition~\ref{prop:dupsize}, and the cofactor size bounds of
Proposition~\ref{prop:seclower} were produced and re-verified by
generator scripts (\lt{codex/gen\_r174.py}, \lt{codex/gen\_r175.py},
\lt{codex/gen\_r197.py}); nothing in them was transcribed by hand. The
generators are not trusted: every identity they emit is re-checked in
Lean by \lt{ring} or \lt{linear\_combination}, and it is the Lean check,
not the generator, that this paper cites.

\section{What is not proved}\label{sec:notproved}

The value of a formalization is the precision of its claims, so we are
explicit.

\begin{enumerate}[label=\textup{(\alph*)}]
  \item \textbf{Lower bounds only.} Theorem~\ref{thm:main} and
        \eqref{eq:389}, \eqref{eq:5077} are rank \emph{lower} bounds.
        Upper bounds require descent --- Selmer groups, the weak
        Mordell--Weil theorem --- which exists in no proof assistant. We
        make no claim about the exact rank of any curve. That the ranks
        of \texttt{389a1} and \texttt{5077a1} are $2$ and $3$ is known
        classically \cite{cremona,lmfdb} and is used nowhere here.
  \item \textbf{No $L$-functions, no BSD.} Nothing here concerns
        $L$-functions, analytic ranks, Tate--Shafarevich groups, or the
        Birch--Swinnerton-Dyer conjecture. No $L$-function of an
        elliptic curve is defined anywhere in the development.
  \item \textbf{The Gram determinant is an input.} Theorem~\ref{thm:main}
        takes $\det\bigl(\langle P_i,P_j\rangle\bigr) \ne 0$ as a
        hypothesis and does not produce it. Certifying it for a given
        curve and a given set of points is a separate, per-curve
        computation --- certified rational interval arithmetic on the
        dyadic chain, as carried out in \cite{MW} for the two curves
        whose certificates are reused here. The pipeline is a machine
        with an input slot, not an oracle.
  \item \textbf{The $2$-torsion hypothesis genuinely restricts the
        class of curves.} Hypothesis (iii) is not removable by better
        bookkeeping: $\Psi$ vanishes exactly at the $x$-coordinates of
        rational $2$-torsion points (\S\ref{sec:ideas}), which is where
        the duplication map and hence the height window break down.
        Curves with rational $2$-torsion --- for instance
        \texttt{102a1}, excluded for the same reason in \cite{MW} ---
        are not covered by any theorem in this paper. A per-step
        variant of the window would be needed, and we have not carried
        it out.
  \item \textbf{Mordell--Weil is neither used nor proved.} All
        statements are about $\rank_{\mathbb{Z}}$ of the point group as
        an abstract $\mathbb{Z}$-module, which is meaningful without
        finite generation. We do not prove that $W(\mathbb{Q})$ is
        finitely generated, and no step depends on it.
  \item \textbf{The two curves are validation, not new arithmetic.}
        \eqref{eq:389} and \eqref{eq:5077} re-derive results first
        obtained in \cite{MW}, and their numerical inputs are the
        certificates of \cite{MW}. They are evidence that the uniform
        machine is non-vacuous and agrees with the hand-built
        development; they are not new facts about elliptic curves.
  \item \textbf{Optimality is not claimed.} The exponent $6$ in
        Theorem~\ref{thm:content}, and the constants of
        \S\ref{sec:price}, are what this route yields. We do not claim
        any of them is best possible, and \S\ref{sec:price} exhibits a
        factor of $\Delta$ that is known to be lost.
\end{enumerate}

\section*{Acknowledgements}

The formalization was carried out in collaboration with Claude
(Anthropic), which wrote and verified the Lean sources under the
author's direction; all mathematical claims were independently rebuilt
and axiom-audited before being recorded. The author thanks the
\texttt{mathlib} community for the elliptic-curve group law, without
which none of this would have been possible.

\begin{thebibliography}{9}

\bibitem{FLT}
K.~Buzzard et al.,
\emph{The Fermat's Last Theorem project},
\url{https://github.com/ImperialCollegeLondon/FLT}, 2024--.

\bibitem{MW}
P.~Cohen,
\emph{Machine-checked canonical heights and independence of rational
points: $\rank \ge 2$ and $\rank \ge 3$ for two elliptic curves},
2026; source at \lt{Papers/formal\_mordell\_weil\_rank\_2026-07-30.tex}
in the repository accompanying this paper.

\bibitem{cremona}
J.~E.~Cremona,
\emph{Algorithms for Modular Elliptic Curves},
2nd ed., Cambridge University Press, 1997.

\bibitem{lmfdb}
The LMFDB Collaboration,
\emph{The $L$-functions and Modular Forms Database},
\url{https://www.lmfdb.org}, 2026.

\bibitem{mathlib}
The mathlib Community,
\emph{The Lean mathematical library},
CPP 2020, 367--381; repository
\url{https://github.com/leanprover-community/mathlib4}.

\bibitem{silverman}
J.~H.~Silverman,
\emph{The Arithmetic of Elliptic Curves},
2nd ed., Graduate Texts in Mathematics 106, Springer, 2009.

\end{thebibliography}

\end{document}
